\section{Proof of the decoder coefficient}\label{sec:proofs}
\subsection{Reducing a repair policy to a random decoder}
\begin{lemma}[Decoder reduction]\label{lem:reduction}
For any ordered policy in $\Pi_\delta$, one can construct, on an enlarged
probability space, its number $N$ of pre-repair observations and a random
decoder $A$ such that
\begin{equation}\label{eq:decodererror}
 \PP_j(A\notin\cD_j)\le\delta\quad\text{for every }j.
\end{equation}
There is auxiliary randomness $U$, whose distribution is independent of the
environment and of the entire pre-repair stream, for which $N$ is a stopping
time in the filtration $\sigma(U,X_1,\ldots,X_n)$ and $A$ is measurable at
$N$. The expected total cost is at least $g_i\EE_iN$.
For a separated policy, $A$ can be taken to be a constant decoder.
\end{lemma}
\begin{proof}
Pre-generate an infinite iid pre-repair stream, including observations unused
by the physical policy. Put policy coins, repair-trial coins, and independent
uniform variables for each possible post-repair class into $U$.
Their joint distribution does not depend on the true environment. Uniform
variables can generate a sample stream from the known law $Q_C$ for each
class $C$.

Follow the policy until it leaves the pre-repair phase, or makes a no-repair
terminal decision. Given the pre-repair history and $U$, its last pre-repair
sample count is a stopping time: no post-repair observation can affect a later
decision to take another pre-repair sample. In the no-repair case, set every
coordinate of $A$ to the final declaration. Otherwise, for each class $C$,
simulate the policy's remaining repair and post-repair continuation using the
same pre-repair internal state, the repair coins, and the synthetic $Q_C$
stream; let $A(C)$ be its final label. A continuation that never succeeds in
repair still produces a deadline declaration under the policy's rule.
The simulations need not be available to the physical policy. They define
an auxiliary statistical answer for the lower bound.

Under environment $j$, the marginal simulated continuation for $C_j$ has
exactly the original policy's continuation law. In particular, its failed
repair trials, any abort decision, and its deadline rule have the correct
distribution. Thus the event $A(C_j)\ne\theta_j$ has the original error
probability, proving \eqref{eq:decodererror}. The other coordinates only
supply a coupling. Both $N$ and $A$ are functions of the pre-repair history
and $U$, and $A$ is measurable at $N$ in the enlarged filtration.
A separated policy's already committed label defines a constant decoder.
Every pre-repair observation costs at least $g_i$, and all other costs are
nonnegative, yielding the cost bound.
\end{proof}

The environment-independent repair mechanism is essential to this reduction.
If success probabilities depend on $i$, their likelihood ratios are extra
information, and the same coefficient need not apply. Likewise, using
unobserved accounting costs as signals would invalidate the experiment.

\subsection{A lower bound with multiple correct answers}
\begin{lemma}[Pre-repair sample lower bound]\label{lem:lower}
If $D_i<\infty$, uniformly over uniformly $\delta$-correct ordered policies,
\[
 \liminf_{\delta\downarrow0}
 \frac{\EE_iN}{L_\delta}\ge\frac{1}{D_i}.
\]
For separated policies the same bound holds with $D_i^{\rm sep}$.
\end{lemma}
\begin{proof}
Fix $i$ and $0<\varepsilon<1$. Every valid decoder has a nonempty invalidating
set when $D_i<\infty$. For each $a\in\cD_i$, choose
$j_a\in B(a)$ attaining $d_i(a)$. Put
\[
 n=\left\lfloor\frac{(1-\varepsilon)L_\delta}{D_i}\right\rfloor,
 \qquad u=(1-\varepsilon/2)L_\delta,\qquad
 L_{ij}(n)=\sum_{t=1}^n\log\frac{P_i(X_t)}{P_j(X_t)}.
\]
The event $F_a=\{N\le n,A=a\}$ is measurable in
$\sigma(U,X_1,\ldots,X_n)$. On this fixed-time sigma-field the density ratio
between environments $i$ and $j_a$ is $\exp L_{ij_a}(n)$, since $U$ has the
same law and is independent of the pre-repair stream. Therefore
\begin{align}
 \PP_i(F_a)
 &\le e^u\PP_{j_a}(F_a)+\PP_i\{L_{ij_a}(n)>u\}\nonumber\\
 &\le \delta^{\varepsilon/2}
       +\PP_i\{L_{ij_a}(n)>u\}. \label{eq:lowerchange}
\end{align}
The second inequality uses $j_a\in B(a)$ and
\eqref{eq:decodererror}. The increments are bounded and iid under $i$, and
\[
 \EE_iL_{ij_a}(n)=n\,d_i(a)\le nD_i\le(1-\varepsilon)L_\delta.
\]
The gap to $u$ is at least $\varepsilon L_\delta/2$.
A bounded-increment concentration bound makes the remaining probability
$o(1)$, independently of the policy. There are finitely many decoders.
Adding \eqref{eq:lowerchange} over valid $a$ and using
$\PP_i(A\notin\cD_i)\le\delta$ gives
$\PP_i(N\le n)=o(1)$. Consequently
$\EE_iN\ge n\PP_i(N>n)=n(1-o(1))$.
First let $\delta\downarrow0$, then $\varepsilon\downarrow0$.

A separated policy permits only constant decoders. The unique constant decoder
valid in $i$ has invalidating set $\{j:\theta_j\ne\theta_i\}$.
Repeating the argument gives the separated rate.
\end{proof}
The proof deliberately uses a fixed-time likelihood ratio on the extended
iid stream. It does not substitute an unstopped likelihood for a stopped one,
and it does not lower-bound each possible answer as if that answer had to be
selected with probability approaching one.

\subsection{A common likelihood-ratio decoder procedure}
At pre-repair sample count $n$, define the stopping rule
\begin{equation}\label{eq:race}
T=\inf\left\{n\ge0:
  \exists a,\ \exists k\notin B(a),\
  L_{kj}(n)\ge\beta_\delta\ \text{for all }j\in B(a)\right\},
\qquad
\beta_\delta=\log\frac{4K(K-1)}{\delta}.
\end{equation}
Empty invalidating sets satisfy the condition at $n=0$.
A deterministic ordering resolves simultaneous eligible decoders and witnesses.
The witness $k$ is searched over all environments; it is not the unknown true
index supplied to the algorithm.

\begin{lemma}[Error and stopping time of the race]\label{lem:race}
The probability that the race ever selects an invalid decoder in environment
$i$ is at most $\delta/4$. If $D_i<\infty$, then
\[
 \limsup_{\delta\downarrow0}\frac{\EE_iT}{L_\delta}\le\frac1{D_i}.
\]
Moreover, for $m=\lfloor H_\delta/3\rfloor$ and sufficiently small $\delta$,
$\PP_i(T>m)\le A_i e^{-a_i m}$ with fixed positive constants $A_i,a_i$.
\end{lemma}
\begin{proof}
Under $i$, an invalid selected decoder must have a witness $k\ne i$
with $L_{ki}(n)\ge\beta_\delta$. The process
$\exp L_{ki}(n)$ is a nonnegative unit-mean likelihood-ratio martingale
under $i$. Ville's maximal inequality and a union bound give
\[
 \PP_i(\text{invalid decoder ever selected})
 \le(K-1)e^{-\beta_\delta}\le\delta/4.
\]
No union over all decoders is needed.

Choose an optimal $a_i\in\cD_i$ and set
$T_i=\inf\{n:\min_{j\in B(a_i)}L_{ij}(n)\ge\beta_\delta\}$.
The race satisfies $T\le T_i$. For a fixed $\eta>0$, let
$n_0=\lceil(1+\eta)\beta_\delta/D_i\rceil$. For $n\ge n_0$,
every $j\in B(a_i)$ has drift $d_{ij}\ge D_i$, and
\[
 nd_{ij}-\beta_\delta\ge nD_i\frac{\eta}{1+\eta}.
\]
If $T_i>n$, at least one likelihood ratio is below threshold at time $n$.
Bounded-increment concentration therefore gives, for some $c_{i,\eta}>0$,
\[
 \PP_i(T>n)\le\PP_i(T_i>n)
 \le |B(a_i)|e^{-c_{i,\eta}n}\quad(n\ge n_0).
\]
Summing tail probabilities bounds $\EE_iT$ by $n_0+o(1)$ as
$\delta\downarrow0$. Divide by $L_\delta$, then let $\eta\downarrow0$.
For the deadline tail use a fixed $\eta$ and
$m/\beta_\delta\to\infty$ to obtain the stated exponential bound.
If a universal decoder exists, $T=0$ and both claims are immediate.
\end{proof}

\subsection{A deadline policy with fully charged tails}
Set $m=\lfloor H_\delta/3\rfloor$. Run the race until it selects a decoder or
until $m$ pre-repair observations have been taken.
If no decoder is selected, finish all remaining slots in the old state
and declare any fixed pre-history rule, such as its maximum-likelihood label.
This event is conservatively counted as error. Otherwise attempt repair for
at most $m$ slots. If all attempts fail, finish the remaining old-state task
slots, charge their costs, and declare a pre-history rule. If repair succeeds,
use the remaining post-repair observations to choose a maximum-likelihood
post class $\widehat C$, and output the selected decoder's label
$a(\widehat C)$. There are at least $m$ post observations after any such success.
For a separated policy, restrict the race to constant decoders and retain
its committed answer even on the all-failed branch.

For distinct post classes $C,D$, let
$\rho_{CD}=\sum_y\sqrt{Q_C(y)Q_D(y)}<1$ and
$\rho=\max_{C\ne D}\rho_{CD}$. Under $Q_C$, for any $r$ observations,
Markov's inequality applied to the square root of the likelihood ratio yields
\[
 \PP_C\!\left\{\prod_{t=1}^r Q_D(Y_t)\ge
                    \prod_{t=1}^r Q_C(Y_t)\right\}
 \le \rho_{CD}^{\,r}.
\]
Thus maximum-likelihood classification among $M=|\calC|$ classes has error at
most $(M-1)\rho^m$ when $r\ge m$.
Take this bound as zero for $M=1$.

The total terminal error is bounded by
\begin{equation}\label{eq:totalerror}
 \frac{\delta}{4}+\PP_i(T>m)+(1-q)^m+(M-1)\rho^m.
\end{equation}
The last three terms are $o(\delta)$, because all instance parameters are
fixed and $m$ grows as $L_\delta^2/3$. Hence the policy belongs to
$\Pi_\delta$ for all sufficiently small $\delta$. For separated policies
the post-class term is unnecessary; the pre-declaration error is already
bounded by $\delta/4+\PP_i(T>m)$.

The expected number of repair attempts on any eligible history is at most
$1/q$. Charging timeout and all-failed tails by the entire possible remaining
old-state cost gives the upper bound
\begin{equation}\label{eq:costupper}
 \EE_i C\le
 g_i\EE_i(T\wedge m)+\frac{\kappa_i}{q}
 +g_iH_\delta\bigl[\PP_i(T>m)+(1-q)^m\bigr].
\end{equation}
This bound includes both failed attempts and unrepaired task slots.
Its last term is $o(1)$. Lemma~\ref{lem:race} gives the desired upper
coefficient. Lemmas~\ref{lem:reduction} and \ref{lem:lower} give the lower
coefficient. Restricting to constant decoders proves the separated equality.
If a universal decoder exists, $T=0$ and \eqref{eq:costupper} is bounded by
$\kappa_i/q+o(1)$. This proves Theorem~\ref{thm:main}.

\begin{remark}[The role of the horizon]
The square-logarithmic deadline is a sufficient explicit schedule. It makes
all exponentially small tail terms negligible relative to $\delta$ and
their task costs negligible relative to $L_\delta$. It is not claimed to be
the smallest feasible deadline. A fixed horizon with bounded likelihood ratios
cannot support arbitrarily small uniform error for all these instances.
\end{remark}

